\documentclass[a4paper,12pt]{article}
\usepackage{fontspec}
\usepackage{polyglossia}
\setotherlanguage{english}
\setmainfont{LiberationSans} % or another font that supports Greek
\usepackage{amsmath, amssymb}
\usepackage[most]{tcolorbox}
\usepackage{xcolor}
\usepackage{colortbl}
\usepackage{enumitem}
\usepackage{multirow}
\usepackage[a4paper, margin=1.8cm]{geometry}
\usepackage{fancyhdr}
\pagestyle{fancy}
\renewcommand{\headrulewidth}{0pt}
\fancyhf{} 
\fancyfoot[L]{\footnotesize Theolaos}
\fancyfoot[R]{\small \thepage}

\usetikzlibrary{decorations.pathreplacing}
\usetikzlibrary{arrows.meta}

\setlist[itemize]{left=1.5em, label=\textbullet}
\setlength{\parskip}{0pt}
\setlength{\parindent}{0pt}

\begin{document}

\begin{center}
    \LARGE \textbf{Σημειώσεις Μαθηματικα Ι}\\[1ex]
    \large Μέθοδοι Ολοκλήρωσεις
\end{center}

\subsection*{Κατά παράγοντες:}
\[
\int f(x)g(x)\,dx = \int F'(x)g(x)\,dx = F(x)g(x) - \int F(x)g'(x)\,dx
\]

\subsection*{Με αντικατάσταση:}
Έστω $f(x) = u$, τότε:
\[
\frac{du}{dx} = u' \iff dx = \frac{du}{u'}
\]
και απλά αντικαθιστώ στη συνέχεια στο ολοκλήρωμα.

ΠΡΟΣΟΧΗ: αντικαταστείς μέσα στο ολοκλήρωμα είτε με τον στόχο για απαλοιφή του $u'$ ή για να αντικαταστήσεις 
στο αποτέλεσμα του $u'$ να το αντικαταστήσεις με $u$ που είχες θέσει. π.χ. 
\[
\int \cos\sqrt{x}\,dx \underset{u = \sqrt{x}}{=} \int \cos{u}\,2u\,du 
\]
που έπειτα το παραπάνω μπορείς να το λύσεις με κατά παράγοντες.

\subsubsection*{Μαθηματικό τρικ (trick):}

Κάποιες φορές στα κλάσματα μπορώ να αντικαταστήσω με σκοπό να βγάλω τη μεταβλητή στον αριθμητή ή και στον παρονομαστή.

\paragraph{Παράδειγμα:}
\[
\int \frac{x}{\sqrt{1-x^2}}\,dx
\]
Έστω $u = 1-x^2$, οπότε $du = -2x\,dx \iff dx = -\frac{du}{2x}$.

Άρα:
\[
\int \frac{x}{\sqrt{u}} \cdot \frac{du}{-2x} = \int \frac{1}{\sqrt{u}} \cdot \frac{du}{-2} = -\int \frac{du}{2\sqrt{u}}
\]
Και φυσικά η λύση θα είναι:
\[
-\sqrt{u} + C
\]

\subsubsection*{Μαθηματικό τρικ (trick):}

Επίσης σαν τρικ μπορείς να χρησιμοποίησεις την παραγοντοποίηση υπερ σου για να πετύχεις τον παραπάνω τρικ.

\paragraph{Παράδειγμα:}
\[
\int \frac{1}{x^2\sqrt{x^2 - 1}}\,dx
\]

Βλέπουμε πως πάνω στον αριθμητή δεν έχουμε κάποιο $x^2$ για να μας σώσει σε κάποια αντικα\-τάσταση
αλλά θα μπορούσαμε να παραγοντοποίησουμε το $x^2-1$ έτσι ώστε να μπορούμε να πάμε ένα $x^a$ κάπως στον αριθμητή ώστε να μας βοηθήσει 
στην μελλοντική αντικατάσταση.

Άρα:
\[
\int \frac{1}{x^2\sqrt{x^2\left(1 - \frac{1}{x^2}\right)}}\,dx = \int \frac{1}{x^3\sqrt{1 - x^{-2}}}\,dx
\]

Εδώ ξεκινάει το τρικ. Άμα θέσουμε $u=x^{-2}$, μπορούμε να φανταστούμε πως η παράγωγος του θα είναι $u' = -2x^{-2 - 1} = -2x^{-3}$.
Άρα άμα προχωρήσω με την αντικατάσταση θα έχω:

\[
\int \frac{1}{x^3\sqrt{1 - u}}\,\frac{du}{-2x^{-3}} = \int \frac{x^3}{-2x^3\sqrt{1 - u}}\,du = \int \frac{1}{-2\sqrt{1 - u}}\,du
\]

Οπότε απλοποιούνται και έχουμε ένα σχετικά απλό ολοκλήρωμα.

\subsection*{Ολοκληρώματα Ρητών Συναρτήσεων}
\hrule
\vspace{0.5em}

Αυτά τα ολοκληρώματα έχουν αυτήν την μορφή $\int \frac{P(x)}{Q(x)}\,dx$ μορφή, όπου $P(x)$ και $Q(x)$ είναι πολυώνυμα.

\

Άμα ο βαθμός του $P(x)$ είναι μεγαλύτερος από αυτόν του $Q(x)$ τότε κάνουμε την πολυώνυμική διαίρεση και θα έχουμε με τον τύπο
της ευκλείδιας διαίρεσης:

\[
\Delta(x) = \delta(x)\cdot\pi(x) + \upsilon(x) \underset{\Delta(x) = P(x)\, \delta(x) = Q(x)}{\text{και άρα}} \int\frac{P(x)}{Q(x)}\,dx = \int\frac{Q(x)\cdot\pi(x) + \upsilon(x)}{Q(x)} = \int \pi(x) + \frac{\upsilon(x)}{Q(x)}\,dx
\]

Όπου το $\pi(x)$ το λύνουμε με την ιδιότητα της γραμμικότητας, και για το κλάσμα κάνουμε τα παρακάτω βήματα:

\begin{itemize}
    \item Βρίσκω τις ρίζες του $Q(x)$
    \item Αναλύω σε περιπτώσεις τις ρίζες του $Q(x)$
    \begin{itemize}
         \item Άμα 
    \end{itemize} 
\end{itemize}

\subsection*{Βασικά Τριγωνομετρικά Ολοκληρώματα}
\hrule
\vspace{0.5em}
\subsubsection*{Περίπτωση μονών δυνάμεων:}
Σε περίπτωση που έχουμε μονές δυνάμεις, π.χ. $\int f(\sin x, \cos x)\,dx$, τότε μπορούμε να εφαρμό\-σουμε την παρακάτω αντικατάσταση:
\[
u = \tan\frac{x}{2}
\]
\[
dx = \frac{2\,du}{1 + u^2}, \quad \sin x = \frac{2u}{1+u^2}, \quad \cos x = \frac{1-u^2}{1+u^2}, \quad \tan x = \frac{2u}{1-u^2}
\]

\subsubsection*{Περίπτωση δυνάμεων εις την δευτέρα:}
Σε περίπτωση που έχεις δύναμη εις την δύο, π.χ. $\int f(\sin^2 x, \cos^2 x)\,dx$, τότε κάνεις την παρακά\-τω αντικατάσταση:
\[
u = \tan x
\]
\[
dx = \frac{du}{1+u^2}, \quad \sin^2 x = \frac{u^2}{1+u^2}, \quad \cos^2 x = \frac{1}{1+u^2}
\]
\subsection*{Υπερβολικά τριγωνομετρικά ολοκληρώματα}
\hrule
\vspace{0.5em}

\subsection*{Υπόλοιπα τριγωνομετρικά ολοκληρώματα}
\hrule
\vspace{0.5em}
\subsubsection*{Ρίζα άθροισμα τετραγώνου ανεξάρτητης μεταβλητής και σταθεράς τιμή:}
Στην περίπτωση που βρεις $\sqrt{(bx)^2 - a^2}$, τότε μπορείς θέσεις τα παρακάτω για να κάνεις αντικα\-τάσταση:

\[
x = \frac{a}{b} \sec \theta,\ \ \ dx = \frac{a}{b}\sec\theta\tan\theta d\theta
\]

Και μετά μπορεις να χρησιμοποίησεις την ταυτότητα:
\[
\sec^2 \theta - 1 = \tan^2\theta
\]

Έπειτα, όταν θες να γυρίσεις πίσω με το $x$, πρέπει να σκεφτείς πως το $\sec\theta = \frac{\text{Υποτείνουσα}}{\text{προσκείμενη}} = \frac{bx}{a}$
άρα μπορούμε να φτιάξουμε κάποιες ταυτότητες αντικατάστασης όπως:
\[
\sin\theta = \frac{b\sqrt{(bx)^2 + a^2}}{bx}
\]
\[
\cos\theta = \frac{a}{bx}
\]
\[
\tan\theta = \frac{\sqrt{(bx)^2 + a^2}}{a}
\]
\[
\theta = \cos^{-1}\left(\frac{a}{bx}\right) \text{ή} \theta = \sec^{-1}\left(\frac{bx}{a}\right) 
\]


\subsection*{Τριγωνομετρικές Ταυτότητες}
\hrule
\vspace{0.5em}
\[
\sin^2 x = \frac{1}{2}(1-\cos 2x)
\]
\[
\cos^2 x = \frac{1}{2}(1+\cos 2x)
\]
\[
\sin x \cos y = \frac{1}{2}(\sin(x-y) + \sin(x+y))
\]
\[
\cos x \cos y = \frac{1}{2}(\cos(x-y) + \cos(x+y))
\]
\[
\sin x \sin y = \frac{1}{2}(\cos(x-y) - \cos(x+y))
\]
\[
\cos(x+y) = \cos x \cos y - \sin x \sin y
\]
\[
\sin(x+y) = \sin x \cos y + \cos x \sin y
\]

\subsection*{Βασικά Αόριστα Ολοκληρώματα}
\hrule
\vspace{0.5em}
Κάποια εξτρά χρήσιμα βασικά αόριστα ολοκληρώματα, πέρα από τα γνωστά, είναι:
\[
\int a^x\,dx = \frac{a^x}{\ln a} + C
\]
\[
\int \frac{1}{\cos^2 x}\,dx = \tan x + C
\]
\[
\int \frac{1}{\sqrt{1-x^2}}\,dx = \arcsin x + C
\]
\[
\int \frac{1}{1+x^2}\,dx = \arctan x + C
\]
\[
\int \frac{1}{x\sqrt{x^2-1}}\,dx = \sec^{-1} x + C
\]

\end{document}
